\clearpage\thispagestyle{empty}
\pdfbookmark[0]{Sources for the September and October 1803 Ceres Report}{ceres18030910-sources}
\begin{center}\textbf{Sources for the September and October 1803 Ceres Report}\end{center}
\noindent This AI translation follows Franz Xaver von Zach’s report as excerpted in \href{https://gdz.sub.uni-goettingen.de/id/PPN236007467}{Gauss’s \textit{Werke}, vol.~VI} (Göttingen, 1874), pp.~245--246. The passage on the revised elements and the two tables correspond to \href{https://api.digitale-sammlungen.de/iiif/image/v2/bsb10538600_00298/full/1450,/0/default.jpg}{p.~290} of article XXXII, \href{https://api.digitale-sammlungen.de/iiif/image/v2/bsb10538600_00296/full/1450,/0/default.jpg}{pp.~288--291}, in the September issue of the \textit{Monatliche Correspondenz}, vol.~VIII (1803). The three observations by Giuseppe Piazzi are printed in the \href{https://api.digitale-sammlungen.de/iiif/image/v2/bsb10538600_00100/full/1450,/0/default.jpg}{July issue, p.~94}. The table in the \textit{Werke} gives Palermo mean times and labels the differences R—B.
\par\smallskip\noindent The bracketed paragraph in the \textit{Werke} describes the ephemeris calculated by Carl Friedrich Gauss and printed in article XLIII, \href{https://api.digitale-sammlungen.de/iiif/image/v2/bsb10538600_00379/full/1450,/0/default.jpg}{pp.~369--371} of the October issue. The ephemeris appears on \href{https://api.digitale-sammlungen.de/iiif/image/v2/bsb10538600_00380/full/1450,/0/default.jpg}{p.~370} and \href{https://api.digitale-sammlungen.de/iiif/image/v2/bsb10538600_00381/full/1450,/0/default.jpg}{p.~371}. Clifford Cunningham gives an English account of the October article in \href{https://doi.org/10.1007/978-3-319-28815-4_10}{“The Scientific Papers on Ceres”} (2016), pp.~364--365, with reproductions of the ephemeris as Figs.~10.68 and~10.69.
\par\smallskip\noindent Translation and typesetting produced by AI. Journal scans from the Bayerische Staatsbibliothek.
\clearpage
